\documentclass[aspectratio=169]{beamer}
\usetheme{Madrid}
\usecolortheme{default}
\usepackage{graphicx}
\usepackage{amsmath}
\usepackage{booktabs}
\usepackage{hyperref}

% UGA Colors
\definecolor{UGARed}{HTML}{BA0C2F}
\definecolor{UGABlack}{HTML}{000000}

\setbeamercolor{palette primary}{bg=UGARed,fg=white}
\setbeamercolor{palette secondary}{bg=UGABlack,fg=white}
\setbeamercolor{palette tertiary}{bg=UGARed,fg=white}
\setbeamercolor{title}{fg=UGARed}
\setbeamercolor{frametitle}{bg=UGARed,fg=white}
\setbeamercolor{structure}{fg=UGARed}

\titlegraphic{\includegraphics[height=1.5cm]{download.png}}
\logo{\includegraphics[height=0.5cm]{download.png}}

\title{When AI Sounds Certain but Is Wrong}
\subtitle{Healthcare, Hallucinations, and Ethical Responsibility}
\author{Breakout 11 Discussion Session}
\date{April 2026}

\begin{document}

% Slide 1
\begin{frame}
\titlepage
\end{frame}

% Slide 2
\begin{frame}[t]{This Week's Theme}
\large
AI hallucination happens when an AI gives information that sounds plausible, confident, and useful, but is actually false or unsupported.

\vspace{0.3cm}
In healthcare, that is not just an inconvenience. A fabricated detail can affect diagnosis, treatment, trust, and patient safety.

\vfill
\textbf{\textcolor{UGARed}{Central Prompt:}}
What makes an AI hallucination especially dangerous in medicine?

\end{frame}

% Slide 3
\begin{frame}[t]{The Med-Gemini Case}
\large
\textit{One reported medical-AI error raised a larger question: typo, or hallucination?}

\vspace{0.3em}

\begin{itemize}
\item Google's Med-Gemini referenced a nonexistent brain structure: ``basilar ganglia.''
\vspace{0.3em}
\item Google suggested it reflected a common mis-transcription learned from training data, but the incident still raises reliability concerns.
\end{itemize}
\vfill
\textbf{\textcolor{UGARed}{Discussion Prompt:}}

If an AI system is built for medical use, is ``it was just a typo'' good enough? What kind of correction or disclosure should the public expect?
\end{frame}

% Slide 4
\begin{frame}[t]{Trust vs. Reliability}
\large
\textit{People may trust AI health information before the evidence fully supports that trust.}

\vspace{0.3em}

\begin{itemize}
\item The article notes that many people already view AI-generated health information as somewhat or very reliable.
\vspace{0.3em}
\item Polished, authoritative language can make weak or false answers sound trustworthy.
\end{itemize}
\vfill
\textbf{\textcolor{UGARed}{Discussion Prompt:}}

What has been your personal experience using AI systems that produce hallucinations or false claims?

\end{frame}

% Slide 5
\begin{frame}[t]{More Powerful, Still Wrong}
\large
\textit{The newer generation of ``reasoning'' models may improve some skills without becoming consistently truthful.}

\vspace{0.3em}

\begin{itemize}
\item The reading reports that hallucinations remain a problem across major AI systems, including newer models marketed as stronger reasoners.
\vspace{0.3em}
\item Better performance in math, coding, or explanation does not automatically make a model more truthful.
\end{itemize}
\vfill
\textbf{\textcolor{UGARed}{Discussion Prompt:}}

Should AI systems be used in high-stakes settings if even their creators do not fully understand why hallucinations persist?
\end{frame}

% Slide 6
\begin{frame}[t]{Healthcare Harm and Bias}
\large
\textit{The problem is not only false facts. It is also unequal harm.}

\vspace{0.3em}

\begin{itemize}
\item A U.S. study described LLMs repeating or elaborating planted false clinical details at troubling rates.
\vspace{0.3em}
\item A U.K. study raised concerns that some models downplayed women's physical and mental health needs in care summaries.
\end{itemize}
\vfill
\textbf{\textcolor{UGARed}{Discussion Prompt:}}

Which is more dangerous in practice: fabricated medical details, biased summaries, or overconfident advice that sounds helpful?
\end{frame}

% Slide 7
\begin{frame}[t]{Who Is Responsible?}
\large
\textit{When AI advice is persuasive but wrong, accountability becomes difficult to hide.}

\begin{itemize}
\item Responsibility can be spread across developers, deployers, clinicians, and institutions.
\vspace{0.3em}
\item ``The human should verify it'' may sound reasonable, but it can also shift risk onto already busy people.
\end{itemize}
\vfill
\textbf{\textcolor{UGARed}{Discussion Prompt:}}

Who should bear the main responsibility when AI advice in a medical setting turns out to be false?
\end{frame}

% Slide 8
\begin{frame}{Closing}
\centering
\textbf{What level of reliability should society require before trusting AI in healthcare?}

\vspace{0.3cm}
After this week's reading, has your view changed on when AI should assist, advise, or stay out of medical decision-making?

\end{frame}

\end{document}
